\section{Introduction}
\label{sec:introduction}

Dementia with Lewy bodies (DLB) is the second most common neurodegenerative dementia after Alzheimer's disease (AD), yet it remains substantially under-recognized in life because its clinical presentation overlaps heavily with AD and Parkinson's disease dementia. The Fourth Consensus Report of the DLB Consortium \cite{mckeith2017dlb} lists reduced occipital metabolism or perfusion on FDG-PET or SPECT as a supportive imaging biomarker, and specifically calls out the ``cingulate island sign'' -- relative metabolic preservation of the posterior cingulate cortex against surrounding hypometabolic parietal and occipital cortex -- as one of the more specific imaging correlates of Lewy body pathology, first characterized quantitatively by Lim et al.\ \cite{lim2009cingulateisland}. These regional metabolic patterns are the biological signal this project's radiomics pipeline is built to characterize at sub-regional, texture level rather than through a single region-averaged uptake ratio.

A persistent obstacle to DLB imaging research is the absence of large, publicly available, pathologically or biomarker-confirmed DLB cohorts. The Alzheimer's Disease Neuroimaging Initiative (ADNI) is the largest public multimodal (PET, MRI, CSF, blood, genetic) longitudinal dementia dataset, but its diagnosis field is restricted to \texttt{CN}, \texttt{MCI}, and \texttt{Dementia} -- it was never designed to label subjects as DLB. ADNI's recent addition of the Amprion CSF alpha-synuclein seed amplification assay (SAA) as a study-wide biomarker changes this: a \texttt{Detected-1} SAA result identifies a synucleinopathy-positive biological state consistent with Lewy body or Parkinson's disease pathology, independent of the clinical diagnosis label attached to the subject. Mak et al.\ \cite{mak2025alphasynuclein} recently used exactly this design -- SAA result, CSF immunoassay, and FDG-PET within a defined time window of one another -- to study the relationship between alpha-synuclein positivity and glucose metabolism across the AD continuum in ADNI, finding SAA-associated hypometabolism concentrated in posterior and occipital cortex. This is the closest published precedent to the present project's cohort-construction strategy, though their analysis is univariate/mediation-based rather than a radiomics classifier.

This project's own cohort was built independently before this literature review, using the same core logic (SAA result as ground truth, FDG-PET and MRI within one year of the SAA draw), yielding 126 SAA-positive and 415 SAA-negative subjects (\texttt{data/adni/dlb\_cohort\_candidates.csv}, \texttt{data/adni/saa\_negative\_controls.csv}). Raw DICOM, ECAT7, and Interfile-format image data for both cohorts, along with ADNI's full tabular data release, biospecimen tables, and derived FreeSurfer and FDG-PET processing outputs, have already been acquired (see \texttt{docs/data\_acquisition.md}). No radiomics pipeline code has yet been written; the present document reviews recent literature to inform five specific pipeline-design decisions that remain open, listed in Section~\ref{sec:methodology}, before implementation begins.

The remainder of this report is structured as follows. Section~\ref{sec:literature_review} summarizes the publications most directly relevant to these design decisions. Section~\ref{sec:proposed_work} states the concrete design choices this review supports. Section~\ref{sec:methodology} frames each of the five open questions as a methodology decision and states the literature-informed answer. Section~\ref{sec:results} summarizes the current implementation status against this design. Section~\ref{sec:conclusion} concludes and lists the residual open items.
